\documentclass[10pt,a4paper]{amsart}
\usepackage[margin=19mm]{geometry}
\usepackage{amsmath,amssymb,mathtools}
\usepackage[scr=rsfs,cal=euler]{mathalfa}
\usepackage{xfrac}
\usepackage[unicode,hidelinks,bookmarks=false]{hyperref}
\newcommand{\ie}{i.e.\ }
\newcommand{\cf}{cf.\ }
\newcommand{\resp}{respectively\ }
\newcommand{\Subsection}{\subsection}
\providecommand{\nobreakdash}{\nobreak\hbox{-}}
\setlength{\emergencystretch}{2em}
\multlinegap0pt
\usepackage[shortlabels]{enumitem}
\usepackage{mathtools}
\usepackage{amssymb}
\usepackage{dsfont}
\usepackage{xcolor}
\usepackage{tikz}
\newtheorem{lemma}{Lemma}
\newtheorem{theorem}{Theorem}
\usepackage{cleveref}
\crefname{lemma}{Lemma}{Lemmas}
\crefname{theorem}{Theorem}{Theorems}
\newcommand{\CC}{\mathbb{C}}
\newcommand{\QQ}{\mathbb{Q}}
\newcommand{\RR}{\mathbb{R}}
\newcommand{\RRt}{\mathbb{R}\{\!\{t\}\!\}}
\DeclareMathOperator{\Trop}{Trop}
\newcommand{\bin}{{\mathrm{bin}}}
\DeclareMathOperator{\grc}{grc}
\DeclareMathOperator{\id}{id}
\DeclareMathOperator{\im}{im}
\newcommand{\lin}{{\mathrm{lin}}}
\DeclareMathOperator{\mrc}{mrc}
\DeclareMathOperator{\rowspan}{rowspan}
\DeclareMathOperator{\val}{val}
\title{Root bounds of vertical systems\\ using tropical~geometry}
\author{Elisenda Feliu, Paul Alexander Helminck, Oskar Henriksson, Yue Ren, Benjamin Schr\"oter, M\'at\'e L. Telek}
\date{Source: arXiv:2605.07645v1 (CC BY 4.0)}
\begin{document}
\maketitle

\noindent\emph{Source and licence.} Root bounds of vertical systems\\ using tropical~geometry. Version arXiv:2605.07645v1 (CC BY 4.0), distributed by arXiv under the Creative Commons Attribution 4.0 International licence, http://creativecommons.org/licenses/by/4.0/, as recorded in the arXiv licence metadata for this exact version and shown on the abstract page https://arxiv.org/abs/2605.07645v1. Full licence notice: \texttt{licenses/arxiv-2605.07645-CC-BY-4.0.txt}.\par\smallskip

\noindent\textit{Editorial scope.} Counted unit: the article's theorem thm:bound\_purevertical (source0.tex lines 1484--1506). The article's complete original proof of it is delivered with it (source0.tex lines 1508--1545, one \texttt{proof} environment of 38 lines). No mathematics is written, repaired or re-derived: the statement and the proof are the article's own lines, in the article's own notation, and no retained line is substituted or re-flowed.  Retained uncounted as the source-local context the proof needs: lines 746--748; lines 1074--1076; lines 1246--1260; lines 1313--1319; lines 1331--1355.  Nothing was omitted from any retained span, so the package carries no omission record.  No retained line contains a citation, so the wrapper carries no bibliography and no bibliography entry.  No retained line contains a figure environment, an image-inclusion command or drawing code, so no image is delivered and the package declares no asset dependency.  Editorial formatting only, in addition: an AMS \texttt{amsart} preamble that restates the article's own declarations and macros used by the retained lines, namely the environments \texttt{lemma}, \texttt{theorem} on separate counters, together with the standard packages those lines need.  The retained environments print in this document as Lemma 1, Theorem 1 in order of appearance, whereas the article prints the same items with the numbering of its own full document.  The complete native source of the article as distributed in the arXiv e-print for this exact version is bundled unchanged as \texttt{sources/200102/source0.tex}.
\begin{equation} \label{eq:monomial_linbin}
I^{\bin}\coloneqq \langle y - x^{M}  \rangle \subseteq \CC[y^\pm,x^\pm]\,, \qquad  I^{\lin} \coloneqq \langle C y  , Lx - b  \rangle \subseteq \CC[a,b][y^\pm,x^\pm]  \,.
\end{equation}

  \begin{equation}\label{eq:shift}
    \Sigma_1\wedge\Sigma_2\coloneqq \lim_{\varepsilon\rightarrow 0} \Sigma_1\wedge (\Sigma_2+\varepsilon\, v)\,,
  \end{equation}

\begin{lemma}\label{lem:technical}
Let $I,J\subseteq \CC[p][x^\pm]$ be prime  complete intersection ideals for parameters $p=(p_1,\dots,p_k)$ 
such that $V(I)$ and $V(J)$ have complementary dimension and $I+J$ has generically nondegenerate zeros. Let $\Trop(I)$ and $\Trop(J)$ be their tropical varieties for generic complex parameter values. 
\begin{enumerate}[label=(\roman*)]
    \item Assume that there exists $\hat{h}\in \QQ^n$ such that $\Trop(I)+\hat{h}$ and $\Trop(J)$ intersect transversely,  and $\Trop(I)+\hat{h}=\Trop(I_p)$ and $\Trop(J)=\Trop(J_p)$ for all $p$ in a  Zariski dense subset $\mathcal{P}$ of $\mathbb{K}^k$.
Then
\begin{equation}\label{eq:technical}
\grc(I+J)= \deg\big((\Trop(I)+\hat{h}) \cap \Trop(J)\big) \, .
\end{equation}
\item Assume that there exists $h\in \QQ^n$ such that for all $\varepsilon>0$ small enough and
$\hat{h} = \varepsilon \, h$,  
$\Trop(I)+\hat{h}$ and $\Trop(J)$ intersect transversely and \eqref{eq:technical} holds. Then
\[ \grc(I+J) = \Trop(I) \cdot \Trop(J)\, . \]
\end{enumerate}
\end{lemma}

\begin{lemma}\label{lem:nu}
    With the notation above, let $(a,b)\in \mathbb{K}^m \times \mathbb{K}^d$ and $z\in \mathbb{K}^r$. Then
    \begin{align*}
    \Trop(I^{\lin}_{a\star \eta(z), b}) & = \Trop(I^{\lin}_{a,b}) + (-\val(z),\mathbf{0}_n)  \\ 
    \Trop^+(I^{\lin}_{a\star \eta(z), b}) & = \Trop^+(I^{\lin}_{a,b}) + (-\val(z),\mathbf{0}_n) \, \qquad  \text{if }z\in \RRt^r_{>0}\,.
    \end{align*}
\end{lemma}

\begin{theorem}[Complex and positive root bounds]
\label{thm:monomial_reembedding_bounds}
Let $F=(C x^M ,\: L x - b)\, \in \CC[a,b][x^\pm]^n$ be a square augmented vertically parametrized system
with   $n=s+d$, $C\in \CC[a]^{s\times r}$ vertical, 
$L\in \mathbb{C}^{d\times n}$ of full row rank, and  $M\in \mathbb{Z}^{n\times r}$. 
Let $I^{\lin}$ and $I^{\bin}$ be as in \eqref{eq:monomial_linbin} and $\mathcal{B}\subseteq \CC^m\times\CC^d$ be the  generic matroid locus of $I^{\lin}$.  

Suppose $h\in\QQ^r$ is such that
$\Trop(I^{\bin})+(h,\mathbf{0}_n)$ and $\Trop(I^{\lin})$ intersect transversely. 
For  $(a^*,b^*)\in\mathcal{B}$, the following statements hold:
\begin{enumerate}[label=(\roman*)]
    \item We have
    \[\grc(F)=\deg\big((\Trop(I^{\bin})+(h,\mathbf{0}_{n} )) \cap \Trop(I^{\lin}_{a^*,b^*})\big)\, .\]

   \item 
  If the coefficients of $F$ are real
    and 
    $(a^*,b^*)\in\RR^m_{>0}\times L(\RR^n_{>0})$, then 
    \begin{align}
\label{eq:num_pos_points1} 
    \mrc_{>0}(F)\geq\#\big( (\Trop^+(I^{\bin})+(h,\mathbf{0}_{n} ))\cap \Trop^+(I_{a^*,b^*}^{\lin})\,\big) \, . 
    \end{align}
   Additionally, the right-hand side of \eqref{eq:num_pos_points1} is a lower bound for the number of nondegenerate positive zeros of  $F_{a\star \eta(t^h),b}$ for all $t>0$ sufficiently small, and generic parameters $(a,b)$ in a  neighbourhood of $(a^*,b^*)$.  
\end{enumerate}
\end{theorem}

\begin{theorem}[Purely vertically parametrized systems]
\label{thm:bound_purevertical}
Let $C \in \mathbb{C}[a]^{n \times r}$ be   a vertical matrix, $a=(a_1,\dots,a_m)$, and $M \in \mathbb{Z}^{n \times r}$ be of rank $n$. Consider the ideals $I^{\lin} \coloneqq \langle\,C y\,\rangle$, with generic matroid locus $\mathcal{B}\subseteq\CC^m$, and $I_M\coloneqq I(\im(\phi_M))$ in $\CC[a][y^\pm]$ in the variables $y=(y_1,\dots,y_r)$. 
\begin{enumerate}[label=(\roman*)]
    \item  It holds that 
\begin{align*}
\grc ( I^{\lin} + I_M  ) &= \rowspan(M)\cdot \Trop(I^{\lin} )\, \\
\grc(C x^M  ) &=  \deg(\phi_M) \, \big(\rowspan(M)\cdot \Trop(I^{\lin})\big)\, .
\end{align*}
\item If $C$ has real coefficients, then 
\begin{align}
\label{eq:max_num_pos_points} \mrc_{>0}(C  x^M ) \leq  \rowspan(M)\cdot \Trop(I^{\lin})  \,, 
\end{align}
and  for  any $a^*\in \mathcal{B}\cap\RR^m_{>0}$
and $h\in\QQ^r$ such that 
$\rowspan(M)+h$ and $\Trop(I^{\lin}_{a^*})$ intersect transversely, it holds that
\begin{align}
\label{eq:num_pos_points} 
\mrc_{>0}(C x^M) \geq \# \big((\rowspan(M)+h)\cap \Trop^+(I^{\lin}_{a^*})\big)\, .
\end{align}
Additionally, the right-hand side of \eqref{eq:num_pos_points} is a lower bound for the number of nondegenerate positive zeros of $C_{a\star \eta(t^h)} x^M$ if $t>0$ is sufficiently small and for generic $a$ in a neighborhood of $a^*$. 
\end{enumerate}
\end{theorem}

\begin{proof} 
Let $F=C x^M$. 
Extending $\phi_M$ to $\mathbb{K}^n$,  we have that
$V( F_a )=\phi_M^{-1}\big( V(C_a y) \cap \im(\phi_M)\big)$ for all  
$a\in \mathbb{K}^m$, 
and thus, $\grc ( F) = \deg(\phi_M)  \grc\big(I^{\lin} + I_M \big)$.
An easy computation gives
\[\Trop(I_M)= \Trop^+(I_M)= \rowspan(M)\,, \]
with a single polyhedron of multiplicity 1. 
Observe that 
$I^{\lin}_a + I_M$ is a complete intersection ideal, and furthermore, 
if $x$ is a nondegenerate zero of $F_a$, then $x^M$ is nondegenerate as a zero of 
$I^{\lin}_a + I_M$. Therefore, 
as $F$ has generically nondegenerate zeros, so does $I^{\lin} + I_M$. 

Let $h\in \QQ^r$ be such that $\Trop(I^{\lin})$ and $\rowspan(M)+ \varepsilon\, h$
intersect transversely  for all $\varepsilon>0$ (see discussion after \eqref{eq:shift}). 
Similar to the proof of \Cref{thm:monomial_reembedding_bounds}, 
consider the automorphism $\mu$ of $\mathbb{K}^m$ defined by $\mu(a)=a\star \eta(t^h)$. Then, as $\mu(\mathcal{B})   \neq \emptyset$ is Zariski dense,  
using \Cref{lem:nu}  and  \Cref{lem:technical}(ii) with the shift $h$, and the ideals $I=I^{\lin}$ and $J=I_M$, we obtain
\begin{equation}\label{eq:toric}
\grc (I^{\lin}+ I_M) =  \rowspan(M)\cdot \Trop(I^{\lin})\, .
\end{equation}
This proves (i). 
To show \eqref{eq:max_num_pos_points} in part (ii), we use that the coefficients are real and that $\phi_M$ is injective when restricted to $\RR^n_{>0}$. Then  for all fixed $a^*\in \RR^m_{>0}$: 
\begin{align*}
&\#\big \{x\in \RR^n_{>0} \mid x  \text{ is a nondegenerate zero of }F_{a^*} \big\}  \\
&\hspace{2cm} =\# \{y\in \RR^r_{>0} \mid y  \text{ is a nondegenerate zero of } I_{a^*}^{\lin} + I_M  \} \\ 
&\hspace{2cm}\leq  \max_{a\in \CC^m} \ \# \{y \in (\CC^*)^r \mid y  \text{ is a nondegenerate zero of } I^{\lin}_{a} +I_M  \} \\
&\hspace{2cm}=\grc ( I^{\lin} + I_M)=\rowspan(M)\cdot \Trop(I^{\lin}) 
\end{align*}
by \eqref{eq:toric}. Equation \eqref{eq:max_num_pos_points}  thus holds.  
The second part of (ii) is simply  \Cref{thm:monomial_reembedding_bounds}(ii) 
 applied to the case where $L$ is empty, 
 as 
\[ \# (\rowspan\,[\,M\mid \id_n\,]+(h,\mathbf{0}_{n} ))\cap \Trop^+(\langle C_a y  \rangle) =\# (\rowspan(M)+h) \cap \Trop^+(\langle C_a y   \rangle)  \, , \]
where the set on the left-hand side belongs to $\RR^{r+n}$. 
\end{proof}

\end{document}
